% Expérience de Millikan : une gouttelette d'huile immobile entre deux plaques (exercice 14, chapitre 11)
\begin{tikzpicture}[x=1cm,y=1cm]
\fill[otMuted!55] (0,2.2) rectangle (5,2.36); \fill[otMuted!55] (0,-0.16) rectangle (5,0);
\fill[otMuted!25] (2.3,2.2) rectangle (2.7,2.36);
\node[pcLegende,anchor=south] at (2.5,2.4) {pulvérisateur d'huile};
\draw[otMuted,line width=0.4pt,{Stealth[length=1.5mm]}-{Stealth[length=1.5mm]}] (5.4,0) -- (5.4,2.2) node[pos=0.5,pcLegende,anchor=west] {$d=\qty{1,0}{cm}$};
\draw[otInk!70,line width=0.5pt] (5.0,2.28) -- (6.6,2.28) -- (6.6,2.7);
\draw[otInk!70,line width=0.5pt] (5.0,-0.08) -- (7.0,-0.08) -- (7.0,2.7);
\node[pcLegende,anchor=south] at (6.8,2.7) {$U=\qty{600}{V}$};
\shade[ball color=yellow!60!orange] (2.5,1.1) circle (0.1);
\node[pcLegende,anchor=east] at (2.32,1.1) {gouttelette};
\draw[otMuted,line width=0.4pt] (-0.4,1.1) -- (0.1,1.1); \node[pcLegende,anchor=east] at (-0.45,1.1) {microscope};
\solution{
  \node[font=\otSemi\fontsize{10}{11}\selectfont,text=otRed!80!black] at (-0.3,2.28) {$+$};
  \node[font=\otSemi\fontsize{10}{11}\selectfont,text=otNavy] at (-0.3,-0.08) {$-$};
  \foreach \xx in {0.3,0.75,4.0,4.6}{\draw[otMuted!80,line width=0.5pt,-{Stealth[length=1.7mm]}] (\xx,2.12) -- (\xx,0.12);}
  \node[pcLegende,anchor=west,fill=white,inner sep=1pt] at (4.65,1.1) {$\vec{E}$};
  \draw[otGreenDark,line width=1.2pt,-{Stealth[length=2.4mm]}] (2.5,1.2) -- (2.5,1.95) node[pcLegende,anchor=west,text=otGreenDark] {$\vec{F}=q\,\vec{E}$};
  \draw[otRed!80!black,line width=1.2pt,-{Stealth[length=2.4mm]}] (2.5,1.0) -- (2.5,0.25) node[pcLegende,anchor=west,text=otRed!80!black] {$\vec{P}=m\,\vec{g}$};}
\end{tikzpicture}
\par\smallskip\legende{On règle la tension pour que la gouttelette observée reste immobile.\solution{ La force électrique compense alors exactement le poids.}}
